The vastness of theoretical drug-like chemical space presents a fundamental challenge for modern drug discovery. Here, we present a generative workflow that addresses this issue through the integration of reinforcement learning-driven molecular generation, adaptive distribution-driven filtering, and transfer learning. The workflow leverages the REINVENT framework and employs adaptive selection criteria based on iteration-specific histogram peaks to avoid reliance on fixed thresholds. Structure-based evaluation is performed using ensemble docking into multiple receptor conformations derived from molecular dynamics simulations, combined with protein-ligand contact surface area (CSA) calculations. As a test case, we applied this approach to the design of novel positive allosteric modulators (PAMs) of the glucagon-like peptide-1 receptor (GLP-1R) acting via a "molecular glue" mechanism to stabilize the endogenous agonist-receptor interface. Across iterations, the workflow yielded consistent enrichment in predicted binding affinity, drug-likeness, and synthetic accessibility. Notably, the top-ranking candidates significantly outperformed the reference compound in docking scores and drug-likeness while successfully recapitulating the glue-like interaction profile. These results demonstrate the ability of the proposed closed-loop framework to steer generative models toward target-relevant regions of chemical space and generate diverse, high-quality candidate molecules.